% !TeX root = ../drafts/05-m3-model.tex
\section{M3 model}\label{sec:m3}

The arithmetization follows the M3 model (multi-multiset matching)~\cite{m3}. The prover chooses each table's size, rows are unordered, and an instance is accepted exactly when every table's constraints hold and the bus balances.

\paragraph{Constraints.} A constraint is a polynomial in the columns that must vanish on every row, including the padding rows.

\paragraph{The bus.} All interactions share a single \emph{bus}, a channel carrying tuples of $m$ field elements. A table flushes by wiring its columns into a tuple and either \emph{pushing} it onto the bus or \emph{pulling} it off, once per row. A few fixed \emph{boundary} tuples, belonging to no table, are pushed or pulled directly. The width $m$ is that of the widest tuple, twelve (the bytecode tuple: separator, $\pc$, count, opcode, and eight operand/immediate slots); shorter tuples are zero-padded.

\paragraph{Domain separation.} The first coordinate of every tuple is a \emph{domain separator}, a fixed constant naming its interaction: $\dsep{ST}$ for state, $\dsep{MEM}$ for memory, $\dsep{BC}$ for bytecode. Since everything shares the one bus, the separators are what keep the three kinds of tuple from cancelling against one another: a memory tuple can never meet a bytecode tuple. The interactions below are three uses of this one channel, told apart by this coordinate.

\paragraph{Balance.} The bus balances when its pushed tuples and its pulled tuples form the same multiset: the same tuples with the same multiplicities, in any order. This is the instance's only global check, proven by grand product argument (\S\ref{sec:gp}).

\paragraph{Per-instruction tables.} Each instruction has its own table, with columns $\pc$, $\fp$, $\mathit{next\_pc}$, $\mathit{next\_fp}$, the operands, the operand addresses, and the read values. For every instruction except \texttt{JUMP}, the successor is virtual: $\mathit{next\_fp}=\fp$ is a copy of $\fp$, and $\mathit{next\_pc}=\gen\cdot\pc$ is the free increment of \S\ref{sec:vm}. Each row makes the same kinds of flush: on the state interaction (\S\ref{sec:state}) it pulls $\tup{\dsep{ST},\pc,\fp}$ and pushes $\tup{\dsep{ST},\mathit{next\_pc},\mathit{next\_fp}}$; it makes one memory flush per cell read (\S\ref{sec:memchan}) and one bytecode flush for the instruction at $\pc$ (\S\ref{sec:bytecode}). The concrete tables are in \S\ref{sec:tables}.

\subsection{Proving the local constraints}\label{sec:air}

Each table must satisfy its own constraints on every one of its own rows. Instead of one zerocheck per table, over cubes of different sizes, we batch them into a single sumcheck.

Table $T_j$ has $2^{\tau_j}$ rows and constraints $C_{j,1},\dots,C_{j,s_j}$: polynomials of degree at most $d$ in a row's column values, required to vanish on every row ($d=2$ throughout). Write $\sigma=\sum_j s_j$ for the total number of constraints and $n=\max_j\tau_j$ for the tallest table's height.

\paragraph{One challenge and one point.} The verifier samples $\eta\in\E$ and fixes a point $r\in\E^{n}$ (in the assembled protocol $r$ is not freshly sampled but taken from the bus, see \emph{Nonzero claims} below and \S\ref{sec:leafstack}). Table $T_j$ mixes its constraints with the powers of $\eta$ from $o_j=s_1+\dots+s_{j-1}$ on, and is tested at the length-$\tau_j$ prefix of $r$:
\[
  C_j\;=\;\sum_{i=1}^{s_j}\eta^{\,o_j+i-1}\,C_{j,i},
  \qquad
  S_j\;=\;\widetilde{C_j}\bigl(r_{<\tau_j}\bigr)\;=\;\sum_{x\in\cube{\tau_j}}C_j(x)\,\eq(r_{<\tau_j},x).
\]
The ranges of $\eta$-powers are \emph{disjoint}, so no table's violation can be cancelled by another's, and proving $\sum_j S_j=0$ proves every $C_{j,i}$ vanishes on every row of its own table, up to Schwartz--Zippel error $(n+\sigma)/|\E|$ over $r$ and $\eta$. (The identities keep disjoint ranges throughout. \S\ref{sec:leafstack} appends three more entries per table that deliberately \emph{share} their powers, since only their total over tables has to be pinned; that lengthens the $\eta$-degree, and the error, by $3$.)

\paragraph{One sumcheck.} The $S_j$ are sumcheck claims of different lengths. Lift each onto the common cube $\cube n$ by attaching the indicator of the all-ones assignment on the $n-\tau_j$ variables $T_j$ does not use:
\[
  f_j(X_0,\dots,X_{n-1})\;=\;\bigl(C_j\,\eq(r_{<\tau_j},\cdot)\bigr)(X_0,\dots,X_{\tau_j-1})\cdot\prod_{i\ge\tau_j}X_i .
\]
Because $\sum_x\prod_i x_i=1$ on any cube, $\sum_{x\in\cube n}f_j(x)=S_j$, so with $F=\sum_j f_j$ the whole statement is the single claim $\sum_{x\in\cube n}F(x)=0$, proven by $n$ rounds of sumcheck on $F$.

The rounds bind $X_{n-1}$ first and $X_0$ last, so the first $n-\tau_j$ touch only variables $T_j$ does not use: there $f_j$ contributes the line $Y\,k\,S_j$, with $k$ the product of the challenges drawn so far (the same for every table that is waiting) and $S_j=0$. Table $T_j$ enters at round $n-\tau_j$ and binds $X_{\tau_j-1},\dots,X_0$ in turn. This is back-loaded batching: a short claim waits, then joins weighted by the rounds it sat out.

\paragraph{Nonzero claims.} Nothing above needs $S_j=0$, nor $r$ to be a fresh challenge. Any polynomial in $T_j$'s columns whose $\eq$-weighted sum over the rows is a value the verifier knows is batched exactly like a constraint, becoming one more entry of $C_j$ with a power of $\eta$ of its own; only $S_j$ changes, from $0$ to that value. And any point fixed at random after the trace commitment serves as $r$. \S\ref{sec:leafstack} does both, running the batch at $r=\zeta$, the point GKR leaves behind, and handing each table three such polynomials. What is proven is then
\[
  \sum_{x\in\cube n}F(x)\;=\;\sum_j S_j .
\]

The one consequence for the rounds is that a waiting table's line $Y\,k\,S_j$ no longer vanishes. Write $X_\ell$ for the variable a round binds, $Y$ for it while still an unknown, and $k$ for the product of the challenges drawn so far, which is common to every waiting table. They therefore sum to $Y\,u$ for the single scalar $u=k\sum_{\tau_j\le\ell}S_j$, and with $p$ the joined tables' cofactor, of degree at most $2$, the round polynomial is the cubic
\[
  h(Y)\;=\;\eq(r_\ell,Y)\,p(Y)+Y\,u .
\]
The prover sends $h$ whole, at $\{0,1,\gen,\gen^{2}\}$: four field elements a round, one more than \S\ref{sec:gkr}'s $\eq$-trick. That one is the price of the nonzero claims rather than of sending $h$ whole, for the trick cannot save it: the verifier has no way to form $u$, so $u$ would travel beside the cofactor instead. The rounds are then vanilla sumcheck, checking $h(0)+h(1)=c$ against the running claim $c$ and folding $c\leftarrow h(\rho_\ell)$, with no $\eq$ reapplied to the message. Prover work is unchanged: $p$ is quadratic, so its fourth value interpolates from the other three rather than costing another pass over the rows.

\subsection{Balancing the bus: grand product}\label{sec:gp}

The bus balances when its pushed and pulled tuples form the same multiset (\S\ref{sec:m3}). Two random challenges, $\alpha$ then $\gamma$, reduce this to a single product.

\emph{Fingerprint.} The verifier samples $\alpha\in\E$ and maps each width-$m$ tuple $\sigma=(\sigma_0,\dots,\sigma_{m-1})$ to a single $\E$-element,
\[
\pi_\alpha(\sigma)\;=\;\sum_{i=0}^{m-1}\alpha^{i}\,\sigma_i ,
\]
each coordinate $\sigma_i$ a committed column value, a virtual column value, or a public constant, all $\K$-elements, so each summand is a mixed $\K\times\E$ product.

\emph{Product check.} The verifier samples $\gamma\in\E$ and the prover proves
\begin{equation}
\prod_{\sigma\ \text{pushed}}\bigl(\gamma-\pi_\alpha(\sigma)\bigr)
\;=\;
\prod_{\sigma\ \text{pulled}}\bigl(\gamma-\pi_\alpha(\sigma)\bigr).
\label{eq:gp}
\end{equation}
\emph{Completeness and soundness.} If the pushed and pulled multisets agree, the two products in \eqref{eq:gp} have the same factors and are equal for every $\alpha,\gamma$; completeness is immediate. For soundness, read each side as the polynomial $\prod_\sigma\bigl(X-\pi_\alpha(\sigma)\bigr)$ evaluated at $X=\gamma$: its roots are that side's fingerprints, which over the formal $\alpha$ are the tuples themselves, so the polynomial determines that side's multiset. The two sides therefore agree as polynomials in $(\alpha,\gamma)$ exactly when the multisets do. When they differ, Schwartz--Zippel applies jointly to the tuple-fingerprinting challenge $\alpha$, the product challenge $\gamma$, the padding correction, and the GKR reductions. If each side has at most $2^\mu$ factors and tuples have width $m$, a conservative total degree bound is
\[
  2m\,2^\mu+8(\mu+1)^2.
\]
The verifier checks that this bound over $|\E|=2^{192}$ leaves at least $128$ bits of soundness. The count product is verified by the same GKR argument, and its root is checked directly for nonzero; it introduces no separate root-at-random test. No Fiat--Shamir grinding is needed for the grand-product phase.

Each side's product is proven by a GKR pass over a product tree (\S\ref{sec:gkr}).

\paragraph{Padding.} Every table is padded past its real rows to a power of two. Each padding row is a fixed default: every column zero but the read counts, set to $\gen^{0}=1$. Its memory and bytecode flushes are then default tuples that do not self-cancel, so the prover announces each table's real row count and the verifier divides the surplus out of \eqref{eq:gp}: a default tuple $d$ in surplus $\delta_d$ on a side divides that side by $(\gamma-\pi_\alpha(d))^{\delta_d}$.


\subsection{Proving the grand product (GKR)}\label{sec:gkr}

Take one side of \eqref{eq:gp}. Its tuples give leaves $v_k=\gamma-\pi_\alpha(\sigma_k)$, padded with $1$s to a power of two, $2^\mu$ in all. Their product $P=\prod_kv_k$ is the root of a binary product tree: layer $0$ is the leaves, layer $i$ has $2^{\mu-i}$ nodes, and layer $\mu$ is $P$. Let $\widetilde V_i:\E^{\mu-i}\to\E$ denote the multilinear extension of layer $i$.

GKR contracts two binary levels at a time. Four grandchildren satisfy
\[
V_i(x)=\prod_{a,b\in\{0,1\}}V_{i-2}(a,b,x),
\]
and hence
\[
\widetilde V_i(r)=\sum_{x\in\cube{\mu-i}}\eq(r,x)\prod_{a,b\in\{0,1\}}\widetilde V_{i-2}(a,b,x).
\]
A sumcheck reduces a claim $\widetilde V_i(r)$ to the four values $\widetilde V_{i-2}(a,b,\rho)$ at one random point $\rho$. The product part of each round polynomial has degree four. As in the usual $\eq$-trick~\cite{Gru24}, the known equality factor is not transmitted. Moreover, the incoming claim already fixes one coefficient. Write the degree-four cofactor as $q(X)=\sum_{j=0}^{4}c_jX^j$, let $s$ be the current coordinate of $r$, and let $T$ be the incoming claim. The prover sends only
\[
d=q(0)+q(1),\qquad c_2,\qquad c_3,\qquad c_4.
\]
Since $T=(1+s)q(0)+sq(1)=c_0+sd$, the verifier recovers $c_0=T+sd$ and $c_1=d+c_2+c_3+c_4$, then evaluates $q$ at the new challenge by Horner's rule. Thus each round sends exactly the four independent field elements permitted by its consistency equation.

At the end of the layer the prover sends the four grandchildren. Two fresh challenges $(u_0,u_1)$ combine them multilinearly, leaving one claim $\widetilde V_{i-2}(u_0,u_1,\rho)$. If $\mu$ is odd, GKR first processes the root-most binary layer, which has no sumcheck rounds, and then proceeds entirely in radix four. After $\lceil\mu/2\rceil$ layers a single leaf claim $\widetilde V_0(\zeta)$ remains, which \S\ref{sec:leafstack} settles against the committed columns. Prover work remains $O(2^\mu)$ and the proof remains $O(\mu^2)$ field elements, with fewer layers and a smaller constant.

\paragraph{Batching the three trees.} The bus needs three grand products: push, pull, and counts. They are batched in a single GKR pass (verifier-friendliness, good for recursion). Push and pull have equal depth $\mu$ (their blocks come in matched pairs); the count tree, of depth $\mu_c\le\mu$, is padded with $1$-leaves to $2^{\mu}$.

The verifier samples a combiner $\lambda$, and each radix-four layer runs one sumcheck on the combined identity
\[
\sum_t\lambda^t\widetilde V_i^t(r)=\sum_{x\in\cube{\mu-i}}\eq(r,x)\sum_t\lambda^t\prod_{a,b\in\{0,1\}}\widetilde V_{i-2}^t(a,b,x),\qquad t\in\{0,1,2\}.
\]
A fresh $\lambda$ after every layer binds the three individual tail values rather than only their previous linear combination. Consequently push, pull, and count all reduce to claims at the same point $\zeta$.

The identity padding is logical: the prover stores only the count tree's actual power-of-two prefix. Once that prefix has folded to one value, each remaining gate pairs it with $1$, while every omitted all-one gate is known and contributes nothing to a nonconstant sumcheck message. The prover also stores each radix-four layer in its original interleaved child order and constructs only the even-numbered binary levels consumed by the protocol. Thus the radix-four verifier savings do not require materializing the padded tree or transposing four child tables.

\subsection{Stacking the leaves}\label{sec:leafstack}

GKR (\S\ref{sec:gkr}) reduces one side of \eqref{eq:gp} to a single evaluation $\widetilde V_0(\zeta)$ of its \emph{leaves}, the product's factors $\gamma-\pi_\alpha(\sigma)$, one per bus tuple $\sigma$. We settle this evaluation against the committed columns.

\paragraph{Flush blocks.} The leaves split into \emph{blocks}, one per flush rule and one per boundary set. A block of a $2^{\kappa}$-row table holds $2^{\kappa}$ leaves; the row-$z$ leaf is $\gamma-\sum_{i=0}^{m-1}\alpha^{i}\,c_{i}(z)$, with $c_{i}(z)$ coordinate $i$ of the row's tuple. Each $c_i$ is public, a committed column, or a virtual column.

\paragraph{Layout.} Order the blocks largest first at aligned offsets, padding to $2^{\mu}$ leaves with $1$s, which leave the product unchanged. This is the stacking of \S\ref{sec:stacking} applied to blocks: block $b$, of $2^{\kappa_b}$ leaves, occupies the subcube $\{(z,\mathsf{sel}_b):z\in\cube{\kappa_b}\}$ for a fixed selector $\mathsf{sel}_b$.

\paragraph{Decomposition.} Split $\zeta=(\zeta^{b}_{\mathrm{lo}},\zeta^{b}_{\mathrm{hi}})$ at block $b$'s boundary. Each block is a subcube, so the leaf MLE splits into a public selector times the block's own MLE, the padding contributing the leftover mass:
\[
  \widetilde V_0(\zeta)\;=\;\sum_b \eq(\mathsf{sel}_b,\zeta^{b}_{\mathrm{hi}})\Bigl(\gamma-\sum_{i=0}^{m-1}\alpha^{i}\,\widetilde c_{b,i}(\zeta^{b}_{\mathrm{lo}})\Bigr)\;+\;\Bigl(1-\sum_b\eq(\mathsf{sel}_b,\zeta^{b}_{\mathrm{hi}})\Bigr).
\]
The verifier forms the selectors $\eq(\mathsf{sel}_b,\cdot)$ and the constant coordinates itself. Each remaining $\widetilde c_{b,i}(\zeta^{b}_{\mathrm{lo}})$ is settled by its kind: a committed column by a direct opening (\S\ref{sec:prelim}), a virtual column off the committed columns it images (so $\gen\cdot\pc$ adds nothing beyond the $\pc$ opening), a \texttt{BLAKE3} value coordinate by a point evaluation of $q_\pkd$ at its slot (\S\ref{sec:tab-blake3}), the index column by \S\ref{sec:idxcol}. A committed column feeding several equal-size blocks is opened once, at their shared point $\zeta_{\mathrm{lo}}$.

\paragraph{Blocks of a table.} Split that sum by owner. Three blocks per side belong to no table, the state boundary and the memory and bytecode seeds (finalizations on the pull side); they settle exactly as above, and the counts side has none of them. Every other block belongs to a table, and a table owns many: one per flush on push and on pull, one per read count on the counts, so \texttt{BLAKE3} owns ten, ten and nine.

All of $T_j$'s blocks are $2^{\tau_j}$ rows tall, so they read its columns at the one point $\zeta_{<\tau_j}$, and that shared point is what collapses them. Fix a side, let $\mathcal B$ be the blocks $T_j$ owns there, and put the row values back where those evaluations stand:
\[
  \sum_{b\in\mathcal B}\eq(\mathsf{sel}_b,\zeta^{b}_{\mathrm{hi}})\Bigl(\gamma-\sum_{i=0}^{m-1}\alpha^{i}\,\widetilde c_{b,i}(\zeta_{<\tau_j})\Bigr)
  \;=\;\sum_{x\in\cube{\tau_j}}\eq(\zeta_{<\tau_j},x)\,B_j(x),
\]
where $B_j$ is that same expression read at a row $x$ of $T_j$, a polynomial every coefficient of which the verifier forms itself (on the counts side $\alpha=1$ and $\gamma=0$, \S\ref{sec:memchan}, so it is a bare sum of count columns). The right-hand side is exactly a claim of \S\ref{sec:air}'s kind, with a nonzero value. So hand the three $B_j$, one per side, to the batch and run it at $r=\zeta$. They take three powers of $\eta$ placed after all $\sigma=\sum_j s_j$ identity powers of \S\ref{sec:air} and \emph{shared} by every table, so with a table's real constraints vanishing,
\[
  S_j\;=\;\eta^{\,\sigma}\,\widetilde B^{\mathrm{push}}_j+\eta^{\,\sigma+1}\,\widetilde B^{\mathrm{pull}}_j+\eta^{\,\sigma+2}\,\widetilde B^{\mathrm{cnt}}_j ,
\]
all at $\zeta_{<\tau_j}$.

Sharing those powers is what ties the batch to the bus. Summing over the tables,
\[
  \sum_j S_j\;=\;\eta^{\,\sigma}R_{\mathrm{push}}+\eta^{\,\sigma+1}R_{\mathrm{pull}}+\eta^{\,\sigma+2}R_{\mathrm{cnt}} ,
\]
each $R_{\mathrm{side}}$ being that side's leaf claim less its table-less blocks and padding, which the verifier forms for itself. The batch therefore starts from a target it derived, and $\eta$ is drawn only afterwards, so reaching that one number forces the three bus equations $\sum_j\widetilde B^{\mathrm{side}}_j=R_{\mathrm{side}}$ at once, and zero in every identity slot besides. Per-table powers would not factor this way: the verifier would need three shares per table transmitted, which is sound but costly, and a transmitted aggregate would settle nothing at all, appearing in no other check.

So the bus sends nothing for a table's blocks, and no column of a table is opened at $\zeta$: no such evaluation travels at all. The batch leaves one claim per column of $T_j$ instead, at its own output point $\rho_{<\tau_j}$, settled as above.

\subsection{State interaction}\label{sec:state}

The state interaction carries $\tup{\dsep{ST},\pc,\fp}$. Each instruction row pulls its current state and pushes its successor $\tup{\dsep{ST},\mathit{next\_pc},\mathit{next\_fp}}$; the bus is preseeded by pushing the initial state $\tup{\dsep{ST},\pc_0,\fp_0}$ and pulling the final state $\tup{\dsep{ST},\pc_{\mathrm{final}},\fp_{\mathrm{final}}}$. When the bus balances, the per-row transitions form a chain from $(\pc_0,\fp_0)$ to $(\pc_{\mathrm{final}},\fp_{\mathrm{final}})$, possibly alongside disjoint cycles. The chain is the execution; a stray cycle is a closed loop of rows that still satisfies every constraint and reads the same committed memory and bytecode, so it cannot affect the start-to-final claim. By convention execution halts at a fixed sentinel: the last bytecode slot. Writing $\nprog$ for the program length (a power of two, \S\ref{sec:e2e-bc}), the instructions sit at $\gen^{0},\dots,\gen^{\,\nprog-1}$, so $\pc_{\mathrm{final}}=\gen^{\,\nprog-1}$, with $\fp_{\mathrm{final}}=\gen^{0}$. A compiled \texttt{main} ends by jumping there. Both boundary states are then public, so the verifier forms them itself.
